\documentclass[12pt]{article}
\usepackage[utf8]{inputenc}
\usepackage[T1]{fontenc}
\usepackage[margin=1in]{geometry}
\usepackage{setspace}
\usepackage{xcolor}
\usepackage[colorlinks=true, urlcolor=blue!60!black, linkcolor=blue!60!black]{hyperref}
\setlength{\parindent}{0em}
\setlength{\parskip}{0.45em}
\renewcommand{\baselinestretch}{1.0}
\begin{document}
\begin{center}
{\Large \textit{Counterman v. Colorado}, 600 U.S. 66 (2023)\par}
\vspace{0.3em}
{\normalsize Case Brief --- One-Page Edition\par}
\end{center}

\textbf{Memory Jogger.} Prosecuting true threats requires proof of at least recklessness, not merely an objective test (\href{https://supreme.justia.com/cases/federal/us/600/22-138/}{600 U.S. 66}).

\textbf{Facts.} Billy Counterman sent hundreds of unwanted Facebook messages to singer Coles Whalen, several of them envisaging violent harm, such as ``Staying in cyber life is going to kill you.'' Terrified, she stopped walking alone and canceled performances.

\textbf{Procedural History.} Convicted of stalking under Colorado's objective true-threats test; the Court of Appeals affirmed.

\textbf{Votes.} Vacated and remanded, 7--2: Kagan for five; Sotomayor (with Gorsuch in part) concurred; Thomas and Barrett dissented.

\textbf{Holding.} The State must prove the speaker ``consciously disregarded a substantial risk that his communications would be viewed as threatening violence.''

\textbf{Reasoning.} The First Amendment gives ``strategic protection'' against chilling speech, as in defamation and incitement; recklessness balances victims' needs and matches \textit{Sullivan}, while incitement, ``a hair's-breadth away'' from political advocacy, still demands intent. Sotomayor would require intent for pure threats; Barrett called recklessness a ``Goldilocks judgment.''

\textbf{Examples.} \emph{Same-Side:} a conviction for a heated livestream comment under a ``reasonable viewer'' instruction. \emph{Opposite-Side:} ominous messages after written notice of fear. \emph{Fence-Sitter:} civil protective orders.

\textbf{Critique.} As Barrett noted, delusional stalkers may be the most dangerous yet the least aware; and the majority's open interest-balancing drew fire from originalists and from speech-protective Justices alike.

\textbf{Key Quotation.} ``We hold that it does, but that a mental state of recklessness is sufficient. The State must show that the defendant consciously disregarded a substantial risk that his communications would be viewed as threatening violence. The State need not prove any more demanding form of subjective intent to threaten another.'' --- Justice Kagan, Opinion of the Court, 600 U.S. at 69.

\end{document}
